% =============================================================================
% Seguimiento del Proyecto: estado en cada avance, cambios y próximos pasos.
% =============================================================================
\section{Seguimiento del Proyecto}
\label{sec:seguimiento}

Este capítulo resume el estado del proyecto al cierre del Avance N°2, los
cambios respecto de lo planteado en el Avance N°1 y el trabajo pendiente.

\subsection{Estado por Avance}

\begin{table}[!htbp]
  \centering
  \caption{Estado de los entregables del proyecto}
  \label{tab:seguimiento}
  \small
  \begin{tabularx}{\linewidth}{@{}>{\raggedright\arraybackslash}p{2.3cm}
                                  >{\raggedright\arraybackslash}p{1.8cm}
                                  >{\raggedright\arraybackslash}X
                                  >{\raggedright\arraybackslash}p{2.2cm}@{}}
    \toprule
    \textbf{Avance} & \textbf{Fecha} & \textbf{Contenido} & \textbf{Estado} \\
    \midrule
    Avance N°1 & 05-11-2026 & Caso de negocio, fuentes de información,
      arquitectura propuesta, modelo estrella y script de creación sin datos
      & Entregado \\
    Avance N°2 & 19-11-2026 & Fuentes descargadas y verificadas; proyecto de
      integración (Python y Apache Hop); esquema de staging; Datawarehouse
      poblado y su script; estadísticas de procesamiento; pruebas de la carga
      & Completado (este informe) \\
    Avance N°3 & 10-12-2026 & Cubo en Cube con dimensiones, métricas y KPI;
      dashboard en Streamlit para la gerencia comercial y de marketing;
      recomendaciones; informe final y presentación & Pendiente \\
    \bottomrule
  \end{tabularx}
\end{table}

\subsection{Cambios Respecto del Avance N°1}

El modelo estrella no cambió: las siete tablas, sus columnas, llaves y reglas
son las del Avance N°1, y el script \bd{01_ddl.sql} se usa sin modificaciones.
Los cambios se concentran en las fuentes y en el reparto de tareas del ETL:

\begin{itemize}[leftmargin=*]
  \item \textbf{Fuente de buses:} el portal de la ANTT (Brasil) rechazó las
        descargas, por lo que la línea de buses usa los datos de pasajeros de
        la Superintendencia de Transporte de Colombia. Como esa fuente no trae
        precio, se agregaron tres fuentes de apoyo para estimar la tarifa y la
        distancia: el SERNAC, el DANE y GeoNames
        (Sección~\ref{sec:tratamientos}).
  \item \textbf{Licencia de Delivery Center:} se encontró la página original
        del dataset; su licencia no es abierta. El grupo decidió usarlo solo
        con fines académicos y no publicar sus archivos ni los datos derivados
        en el repositorio.
  \item \textbf{Asignación de clientes:} se hace en la preparación en Python, y
        no en Apache Hop como se planteó en el Avance N°1, porque requiere
        muestreo ponderado. Hop conserva la conversión de monedas y unidades,
        la traducción de códigos, las reglas de calidad y la carga.
  \item \textbf{Taxis:} se usan los taxis verdes de Nueva York, incluidos en la
        fuente aprobada, en vez de los amarillos: sus archivos pesan unos
        1,3~MB al mes, frente a 50 a 70~MB, y la muestra es la misma.
  \item \textbf{Aerolínea:} BTS todavía no publica el segundo semestre de
        2025, que se completa con 2024 trasladado un año.
\end{itemize}

\subsection{Próximos Pasos}

Para el Avance N°3 quedan pendientes:

\begin{enumerate}[leftmargin=*]
  \item Levantar Cube en Docker Compose y definir el cubo con solo las
        dimensiones y atributos necesarios para el perfil táctico, y las
        métricas de la Sección~\ref{sec:kpi} (ventas netas, ticket promedio,
        clientes activos, frecuencia de compra, gasto anual por cliente,
        participación por línea, clientes multilínea y tasa de descuento).
  \item Construir el dashboard en Streamlit y Plotly, que responda las
        preguntas de negocio de la Sección~\ref{sec:preguntas}.
  \item Redactar las recomendaciones de negocio y la presentación.
\end{enumerate}
